\documentclass[border=10pt]{standalone}
\usepackage{tikz,ctex}
\usetikzlibrary{positioning}

\begin{document}

\begin{tikzpicture}[
    % 全局设置
    node distance = 1.6cm and 1.4cm,
    every node/.style = {
        draw,
        rounded corners,
        thick,
        align = center,
        font = \sffamily\small,
        minimum width = 2.0cm,
        minimum height = 0.9cm
    },
    % 样式定义
    centerStyle/.style = {
        fill = blue!45,
        text = white,
        font = \sffamily\bfseries\large,
        line width = 2pt,
        minimum size = 2.8cm,
        circle
    },
    branchTitle/.style = {
        font = \sffamily\bfseries,
        line width = 1.5pt
    },
    subNode/.style = {
        fill = white,
        font = \sffamily\small,
        minimum width = 2.2cm,
        minimum height = 0.8cm
    },
    leafNode/.style = {
        fill = white,
        font = \sffamily\footnotesize,
        minimum width = 1.9cm,
        minimum height = 0.7cm
    },
    arrowStyle/.style = {
        ->,
        thick,
        >=stealth,
        shorten >=2pt
    }
]

    % ================= 1. 中心节点 =================
    \node[centerStyle] (center) {多层网络 \\ Multilayer Networks};

    % ================= 2. 上方分支：网络结构与参数 (红) =================
    \node[branchTitle, fill=red!15, above=of center] (structure) {网络结构与参数 \\ Network Structure \& Parameters};

    \node[subNode, above=of structure, xshift=-6cm] (twolayer) {两层网络 \\ Two-layer Network};
    \node[subNode, above=of structure, xshift=-2cm] (preact) {预激活 \\ Pre-activation $a_j^{(1)}$};
    \node[subNode, above=of structure, xshift=2cm] (hidden) {隐藏单元 \\ Hidden Units $z_j^{(1)}$};
    \node[subNode, above=of structure, xshift=6cm] (output) {输出单元 \\ Output Units $y_k$};

    % ================= 3. 右侧分支：参数矩阵表示 (蓝) =================
    \node[branchTitle, fill=orange!15, right=of center] (matrix) {参数矩阵表示 \\ Parameter Matrices};

    \node[subNode, right=of matrix, yshift=4cm] (biasabsorb) {偏置吸收 \\ Bias Absorption $x_0=1$};
    \node[subNode, right=of matrix, yshift=2cm] (overall) {整体函数 \\ Overall Function (6.11)};
    \node[subNode, right=of matrix, yshift=0cm] (vector) {向量形式 \\ Vector Form (6.12)};
    \node[subNode, right=of matrix, yshift=-2cm] (weightmatrix) {权重矩阵 \\ Weight Matrices $\mathbf{W}^{(1)}, \mathbf{W}^{(2)}$};
    \node[subNode, right=of matrix, yshift=-4cm] (elementwise) {逐元素激活 \\ Element-wise Activation};

    % ================= 4. 下方分支：万能逼近与激活函数 (绿) =================
    \node[branchTitle, fill=green!15, below=of center] (approx) {万能逼近与激活函数 \\ Universal Approximation \& Activations};

    \node[subNode, below=of approx, xshift=-10cm] (universal) {万能逼近定理 \\ Universal Approximation};
    \node[subNode, below=of approx, xshift=-5cm] (limitation) {局限性 \\ Exponential Units, No Free Lunch};
    \node[subNode, below=of approx, xshift=0cm] (sigmoid) {Sigmoid / Tanh \\ 梯度消失};
    \node[subNode, below=of approx, xshift=3.5cm] (relu) {ReLU / Softplus \\ 非零梯度};
    \node[subNode, below=of approx, xshift=7cm] (leaky) {Leaky ReLU \\ 可学习 $\alpha_j$};

    % ================= 5. 左侧分支：权重空间对称性 (紫) =================
    \node[branchTitle, fill=purple!15, left=of center] (symmetry) {权重空间对称性 \\ Weight-space Symmetries};

    \node[subNode, left=of symmetry, yshift=4cm] (signflip) {符号翻转对称 \\ Sign-flip Symmetry $2^M$};
    \node[subNode, left=of symmetry, yshift=2cm] (interchange) {交换对称 \\ Interchange Symmetry $M!$};
    \node[subNode, left=of symmetry, yshift=0cm] (factor) {总对称因子 \\ Overall Factor $M!2^M$};
    \node[subNode, left=of symmetry, yshift=-2cm] (multilayer) {多层推广 \\ Product over Layers};
    \node[subNode, left=of symmetry, yshift=-4cm] (bayesian) {贝叶斯方法中的作用 \\ Role in Bayesian Methods};

    % ================= 连线逻辑 =================
    % 中心 -> 四大分支标题
    \draw[arrowStyle, red] (center) -- (structure);
    \draw[arrowStyle, blue] (center) -- (matrix);
    \draw[arrowStyle, green!60!black] (center) -- (approx);
    \draw[arrowStyle, purple] (center) -- (symmetry);

    % 分支标题 -> 子节点 (上方)
    \draw[arrowStyle, red!60] (structure) -- (twolayer);
    \draw[arrowStyle, red!60] (structure) -- (preact);
    \draw[arrowStyle, red!60] (structure) -- (hidden);
    \draw[arrowStyle, red!60] (structure) -- (output);

    % 分支标题 -> 子节点 (右侧)
    \draw[arrowStyle, blue!60] (matrix) -- (biasabsorb);
    \draw[arrowStyle, blue!60] (matrix) -- (overall);
    \draw[arrowStyle, blue!60] (matrix) -- (vector);
    \draw[arrowStyle, blue!60] (matrix) -- (weightmatrix);
    \draw[arrowStyle, blue!60] (matrix) -- (elementwise);

    % 分支标题 -> 子节点 (下方)
    \draw[arrowStyle, green!60!black] (approx) -- (universal);
    \draw[arrowStyle, green!60!black] (approx) -- (limitation);
    \draw[arrowStyle, green!60!black] (approx) -- (sigmoid);
    \draw[arrowStyle, green!60!black] (approx) -- (relu);
    \draw[arrowStyle, green!60!black] (approx) -- (leaky);

    % 分支标题 -> 子节点 (左侧)
    \draw[arrowStyle, purple!60] (symmetry) -- (signflip);
    \draw[arrowStyle, purple!60] (symmetry) -- (interchange);
    \draw[arrowStyle, purple!60] (symmetry) -- (factor);
    \draw[arrowStyle, purple!60] (symmetry) -- (multilayer);
    \draw[arrowStyle, purple!60] (symmetry) -- (bayesian);

\end{tikzpicture}

\end{document}